\chapter{Free fields on an interval}
\label{chap:free-field-interval}

A free action determines Gaussian contractions.  It does not determine a
multiplication on a chosen complex of topological cochains.  The distinction
already appears when a field has values at two endpoints and a descendant
integrated along the interval between them.  Linear interpolation gives a
binary product which is associative only up to a specified ternary operation.
For a Heisenberg current, that operation retains the double pole of the
operator product expansion.  Its coefficient is an elementary integral.

The construction below keeps the holomorphic insertion points distinct.
It gives all ordered higher operations on the polynomial de Rham model of
an interval, with coefficients in the multilocal Wick algebra.  An explicit
comparison identifies the Wick factors with radially ordered free fields.
Topological cochain degree, holomorphic jet order, and powers of the quantum
parameter remain separate throughout.

\section{The current and its separated products}

Work on the affine complex curve with coordinate $z$.  Put
$K=\C[\kappa]\llbracket\hbar\rrbracket$, where both parameters have
cohomological degree zero.  The rank-one Heisenberg vacuum module over $K$ is
\[
 V=\C[\kappa][q_1,q_2,\ldots]\llbracket\hbar\rrbracket.
\]
Each coefficient of $\hbar$ is a polynomial in finitely many $q_m$.
All vacuum generators and holomorphic jet symbols have cohomological degree zero.
For $m>0$, define operators
\[
 J_{-m}=q_m,\qquad
 J_m=\hbar\kappa m\frac{\partial}{\partial q_m},\qquad J_0=0.
\]
The vacuum is $\mathbf1=1$.  Direct differentiation gives
\[
 [J_m,J_n]=\hbar\kappa m\delta_{m+n,0}\,\id.
\]
The field $J(z)=\sum_{n\in\Z}J_nz^{-n-1}$ splits as
$J_-(z)+J_+(z)$, with creation modes in $J_-$ and annihilation modes
in $J_+$.  In the formal expansion $|z|>|w|$,
\begin{equation*}
 [J_+(z),J_-(w)]
 =\hbar\kappa\sum_{m\ge1}m z^{-m-1}w^{m-1}
 =\frac{\hbar\kappa}{(z-w)^2}.
\end{equation*}
Here and below, a rational expression denotes its specified formal Laurent
expansion.  No norm convergence of unbounded operators is asserted.

Normal ordering places every creation mode to the left of every
annihilation mode.  The field attached to
$q_{r_1+1}\cdots q_{r_a+1}\mathbf1$ is
\[
 :\prod_{\nu=1}^a\frac{\partial_z^{r_\nu}J(z)}{r_\nu!}:.
\]
The denominators identify mode indices with divided holomorphic jets.
All holomorphic derivatives are retained.

For a finite set $S$ of insertion labels, define
\[
 R_S=\C[z_i\ (i\in S),\ (z_i-z_j)^{-1}\ (i\ne j)],
 \qquad
 W_S=\bigl(R_S[\kappa][x_{i,r}\mid i\in S,\ r\ge0]\bigr)
            \llbracket\hbar\rrbracket.
\]
Tensor products of field modules are over $K$, completed in $\hbar$,
with the displayed extension of rational scalars on the output.
The symbol $x_{i,r}$ represents the undivided derivative
$\partial_{z_i}^rJ(z_i)$ inside one globally normal-ordered expression.
Every coefficient is a polynomial in finitely many jet symbols.  For disjoint
sets $S,T$, extend scalars to $R_{S\cup T}$ and set
\begin{equation*}
 \begin{split}
 P_{S,T}&=\sum_{i\in S,\ j\in T}\sum_{r,s\ge0}
 C_{ij}^{rs}\frac{\partial}{\partial x_{i,r}}
              \frac{\partial}{\partial x_{j,s}},\\
 C_{ij}^{rs}&=\hbar\kappa\partial_{z_i}^r\partial_{z_j}^s
                    (z_i-z_j)^{-2},\qquad
 f\star g=\exp(P_{S,T})(fg).
 \end{split}
\end{equation*}
The derivatives in $P_{S,T}$ act only on jet symbols.  They do not act on
its rational coefficients.  On polynomial inputs the exponential is finite.
For formal inputs it is defined coefficientwise in $\hbar$.

Fix a total order $i_1,\ldots,i_a$ on $S$, with
$|z_{i_1}|>\cdots>|z_{i_a}|$ as a formal expansion convention.
Writing $V_0=\C[\kappa][q_1,q_2,\ldots]$, define
\[
 \mathcal N_S:W_S\longrightarrow
 \Hom_{\C[\kappa]}\bigl(V_0,
 V_0((z_{i_1}))\cdots((z_{i_a}))\bigr)\llbracket\hbar\rrbracket
\]
by replacing each jet monomial with its globally normal-ordered field
and expanding each rational coefficient in that order.  A fixed vacuum
polynomial involves only finitely many annihilation modes, so the indicated
iterated Laurent series are defined coefficientwise.

\begin{proposition}[Associativity and the field comparison]
\label{prop:interval-wick-associativity}
For pairwise disjoint $S,T,U$, the two maps
\[
 W_S\otimes W_T\otimes W_U\longrightarrow W_{S\cup T\cup U},
 \qquad (f\star g)\star h,\quad f\star(g\star h)
\]
are equal.  In a radial ordering with every point of $S$ outside every
point of $T$, the normal-ordering map $\mathcal N$ satisfies
\[
 \mathcal N(f\star g)=\mathcal N(f)\mathcal N(g).
\]
Both identities hold for every holomorphic jet and every power of $\hbar$.
\end{proposition}
\begin{proof}
All contraction operators commute: their coefficients are independent of
the jet symbols being differentiated.  Moreover,
$P_{S\cup T,U}=P_{S,U}+P_{T,U}$.  Each parenthesization is therefore
\[
 \exp(P_{S,T}+P_{S,U}+P_{T,U})(fgh).
\]
This proves associativity without identifying distinct insertion points.

For the comparison, commute each annihilation mode in $f$ past each
creation mode in $g$.  The displayed annihilation--creation commutator
gives the contribution of one commutation.  Its derivatives give
$C_{ij}^{rs}$.  Every commutator is central, so further commutations do not
alter it.  A collection of $a$ disjoint contractions occurs $a!$ times in
$P_{S,T}^a$; division by $a!$ counts it once.  Uncontracted modes remain
in normal order.  This proves the formula on monomials, hence on
polynomials.  Formal extension is coefficientwise: each input coefficient
has finite polynomial degree, and only finitely many powers of $\hbar$
contribute to a fixed coefficient.
\end{proof}

For each fixed radial order, put $W_S^{\mathrm{fld}}=\operatorname{im}\mathcal N_S$.
The displayed comparison identifies its product with composition of the
actual field operators on ordered blocks.  If $\mathcal N_S(f)=0$, then
$\mathcal N_{S\cup T}(f\star g)=0$ for every following block $T$.
The same holds for a zero factor in the second position.
Thus the kernels are preserved by all ordered products.  Every interval
operation constructed below is a scalar cochain operation multiplied by
such a product.  It therefore descends to
$C_I\otimes W_S^{\mathrm{fld}}$, with the same contraction and proofs.
This quotient supplies the operations on the actual free fields, including
any relations among their normally ordered expressions.

Write $x_i=x_{i,0}$ and $C_{ij}=\hbar\kappa/(z_i-z_j)^2$.
The first products are
\begin{equation*}
 \begin{split}
 x_1\star x_2&=x_1x_2+C_{12},\\
 x_1\star x_2\star x_3
 &=x_1x_2x_3+C_{12}x_3+C_{13}x_2+C_{23}x_1.
 \end{split}
\end{equation*}
Vacuum expectation means taking the scalar term after normal ordering.
Two copies of $J$ at the first insertion give
\[
 \langle :J^2:(z_1)J(z_2)J(z_3)\rangle_{\mathrm{conn}}
   =2C_{12}C_{13}.
\]
For three quadratic insertions the connected expectation is
$8C_{12}C_{13}C_{23}$.  There are two assignments of the first pair
of fields to the other two vertices.  At each remaining vertex there
are two choices of the field used for its first edge.  This gives
$2\cdot2\cdot2=8$ triangle pairings.  Thus a Gaussian theory has
nonzero connected composite-field products.

\section{Endpoint values and an integrated descendant}

Give the interval $I=[0,1]$ its increasing orientation.  Its polynomial
forms are the differential graded algebra
\[
 A_I=\C[t,dt],\qquad |t|=0,\quad |dt|=1,\quad d t=dt,
 \qquad (dt)^2=0.
\]
The complex of cellular cochains is
\[
 C_I^0=\C e_0\oplus\C e_1,\qquad C_I^1=\C\eta,
 \qquad de_0=-\eta,\quad de_1=\eta,\quad d\eta=0.
\]
Here $e_0,e_1$ record endpoint values and $\eta$ records integration
along the oriented edge.  There are no endpoint vanishing conditions.
Define inclusion and projection by
\begin{equation*}
 \begin{aligned}
 i(e_0)&=1-t,\qquad i(e_1)=t,\qquad i(\eta)=dt,\\
 p(f)&=f(0)e_0+f(1)e_1,\\
 p(g\,dt)&=\left(\int_0^1g(t)\,dt\right)\eta.
 \end{aligned}
\end{equation*}
The degree-minus-one homotopy is
\begin{equation*}
 h(f)=0,\qquad
 h(g\,dt)=\int_0^t g(u)\,du-t\int_0^1g(u)\,du.
\end{equation*}
It subtracts the linear function with the same endpoint difference.

\begin{lemma}[The interval contraction]
\label{lem:interval-contraction}
The maps above satisfy
\[
 pi=1,\qquad dh+hd=1-ip,\qquad hi=ph=h^2=0.
\]
\end{lemma}
\begin{proof}
The inclusion and projection commute with $d$ by the fundamental theorem
of calculus.  On $f(t)$, the operator $hd$ gives
$f(t)-f(0)-t(f(1)-f(0))$.  On $g(t)dt$, the operator $dh$ gives
$(g(t)-\int_0^1g)dt$.  These are exactly $1-ip$ in their respective
degrees.  The function $h(gdt)$ vanishes at both endpoints, proving
$ph=0$.  Also $h(dt)=0$, $h$ vanishes on functions, and its image
consists of functions.  These observations give $hi=h^2=0$.
\end{proof}

Use the complexes
\[
 A(S)=A_I\otimes W_S,\qquad C(S)=C_I\otimes W_S,
\]
completed coefficientwise in $\hbar$.  Their differential acts on the
interval factor only.  For disjoint supports their product on $A$ is
\[
 (\alpha f)(\beta g)=(\alpha\wedge\beta)(f\star g).
\]
Proposition~\ref{prop:interval-wick-associativity} makes this an
associative system of differential graded products.  The maps $i,p,h$
act on its interval factor and preserve each holomorphic support.
These are topological forms with chiral field coefficients.  This
specified algebraic model does not identify the interval differential
with a quantum Batalin--Vilkovisky differential.

The projected binary product $m_2=p\mu(i\otimes i)$ obeys
\[
 \begin{gathered}
 m_2(e_a,e_b)=\delta_{ab}e_a,\qquad
 m_2(e_a,\eta)=m_2(\eta,e_a)=\tfrac12\eta,\\
 m_2(\eta,\eta)=0 \qquad (a,b\in\{0,1\}).
 \end{gathered}
\]
Every entry is multiplied by the corresponding Wick product of its
field coefficients.  In particular,
\[
 m_2(m_2(e_1,e_1),\eta)-m_2(e_1,m_2(e_1,\eta))
 =\tfrac14\eta.
\]
Linear interpolation therefore forces a higher homotopy.

\section{The higher operations and their signs}

For a cochain $a$, write $sa$ for its suspension, with
$|sa|=|a|-1$.  A sequence of ordered operations is encoded by
maps $b_n$ of degree one on suspended inputs.  For pairwise disjoint
supports these satisfy
\begin{equation*}
 \sum_{r+q+t=n}(-1)^{|u_1|+\cdots+|u_r|}
 b_{r+1+t}(u_1,\ldots,u_r,b_q(u_{r+1},\ldots,u_{r+q}),
                         u_{r+q+1},\ldots,u_n)=0.
\end{equation*}
All exponents use suspended degrees.  The ordinary operations $m_n$
of degree $2-n$ are defined by
\begin{equation*}
 b_n(sa_1,\ldots,sa_n)
 =(-1)^{\sum_{j=1}^n(n-j)|a_j|}\,s m_n(a_1,\ldots,a_n).
\end{equation*}
Here the tensor coalgebra has a precise support convention.  Fix a
finite ambient label set $E$ and extend rational scalars to $R_E$.
Take the direct sum of tensor words whose entries have pairwise disjoint
supports $S_1,\ldots,S_n\subseteq E$.  The coproduct cuts a word into
two consecutive pieces.  Each operation replaces a consecutive block by
one entry supported on its union.  Thus all displayed compositions are
defined over a common coefficient ring and remain in this coalgebra.
The formulas commute with further extensions of $E$.  An $A_\infty$
structure here means a square-zero coderivation of these supported words.
The words are finite, so no sum over arities is needed on an element.

On $sA$, put $D_1(sa)=s(da)$ and
$D_2(sa,sb)=(-1)^{|a|}s(ab)$.  The differential graded algebra
identities say $D^2=0$, where $D=D_1+D_2$ is the induced coderivation.
Write $I,P,H$ for the suspended maps $i,p,h$.  Starting with
$F_1=I$, define recursively, for $n\ge2$,
\begin{equation*}
 R_n=\sum_{a+b=n}D_2(F_a\otimes F_b),\qquad
 F_n=-H R_n,\qquad b_n=P R_n,
 \qquad b_1=s d s^{-1}.
\end{equation*}
Each $F_n$ has degree zero.  Thus its tensor substitution contributes
no extra sign.  The only signs come from $D_2$ and the coderivation
insertion rule in the displayed higher identity.

\begin{theorem}[Ordered transfer with free-field coefficients]
\label{thm:interval-free-field-transfer}
The recursion defining $R_n,F_n,b_n$ gives an
uncurved $A_\infty$ structure on the collection $\{C(S)\}$, with operations for every
ordered list of pairwise disjoint supports.  The maps $F_n$ define
an $A_\infty$ quasi-isomorphism to $\{A(S)\}$.  All identities preserve
the rational holomorphic coefficients and their formal $\hbar$ expansion.
The element $e_0+e_1$ with empty field support is a strict unit.
\end{theorem}
\begin{proof}
The proof takes place at each finite word length.  Suppose the
morphism equation $DF=Fb$ and $b^2=0$ hold in all smaller lengths.
Set the length-$n$ components temporarily to zero.  The length-$n$
morphism defect is
\[
 Q_n=R_n-L_n,\qquad
 L_n=\sum_{\substack{r+q+t=n\\2\le q\le n-1}}
 F_{r+1+t}(1^{\otimes r}\otimes b_q\otimes1^{\otimes t}),
\]
with the tensor signs fixed above.  Let $O_n$ be the length-$n$
component of $b^2$ with $b_n=0$.  The identity
\[
 D(DF-Fb)+(DF-Fb)b=-Fb^2
\]
gives $\delta Q_n=-I O_n$, where
$\delta Q_n=D_1Q_n+Q_n b_1$ and $b_1$ on the right denotes its
extension to the input word.  Lower defects vanish by induction.
For $k\ge2$ the recursion gives $PF_k=HF_k=0$.  Consequently
$PL_n=HL_n=0$.

Set $b_n=PQ_n=PR_n$.  Applying $P$ to the defect identity gives
$\delta b_n=-O_n$, which is exactly the length-$n$ equation $b^2=0$.
Next set $F_n=-HQ_n=-HR_n$.  The contraction identity and $HI=0$
give
\[
 \delta F_n=-Q_n+IPQ_n+H\delta Q_n=I b_n-Q_n.
\]
This is the missing length-$n$ morphism equation.  Induction proves
both assertions.  The linear term $I$ is a quasi-isomorphism by
Lemma~\ref{lem:interval-contraction}.

For the unit assertion, consider a binary transfer tree with a unit
leaf.  At its first vertex the unit removes that multiplication.
If the other branch is a leaf, the next homotopy kills it by $HI=0$.
If that branch has an internal vertex, the next map gives $H^2=0$
or $PH=0$.  Thus every operation of arity greater than two vanishes
on a unit input.  The binary and unary unit identities follow from
the displayed table and $d(e_0+e_1)=0$.

Every tree has finitely many vertices.  Associativity makes its field
coefficient the same ordered Wick product, independently of the tree.
The polynomial contraction acts over $K$ and neither changes a holomorphic denominator
nor a power of $\hbar$.  The construction therefore extends to the
stated coefficients and all holomorphic jets.
\end{proof}

The Wick exponential supplies every field coefficient of the transfer.
The explicit interval contraction fixes the higher operations,
including choices which a boundary identity alone cannot fix.

\section{The ternary operation and the next identities}

Write $v=se_1$, $w=s\eta$, and $\mathbf e=s(e_0+e_1)$.
Their degrees are $-1,0,-1$, respectively.  The first suspended
operations, with field coefficients suppressed, are
\[
 \begin{gathered}
 b_1(v)=w,\qquad b_1(w)=0,\qquad
 b_2(v,v)=v,\qquad b_2(v,w)=\tfrac12w,
 \qquad b_2(w,v)=-\tfrac12w,\\
 b_3(v,w,w)=\tfrac1{12}w,\qquad
 b_3(w,v,w)=-\tfrac16w,\qquad
 b_3(w,w,v)=\tfrac1{12}w.
 \end{gathered}
\]
Every other ternary value on $v,w$ is zero.  Values involving $se_0$
follow from $se_0=\mathbf e-v$ and strict unitality.

The integral producing the first coefficient is
\[
 h(t\,dt)=\frac{t^2-t}{2},\qquad
 F_2(v,w)=s\frac{t-t^2}{2},\qquad
 b_3(v,w,w)=s\left(\int_0^1\frac{t-t^2}{2}\,dt\right)\eta
 =\frac1{12}w.
\]
For the middle input there are two trees, each with the opposite sign.
This gives $-w/6$.  The last input gives the same value as the first.

\begin{example}[The Heisenberg ternary operation]
\label{ex:interval-heisenberg-m3}
For three distinct holomorphic insertions, ordinary unsuspension gives
\begin{equation*}
 \begin{split}
 &m_3(e_1x_1,\eta x_2,\eta x_3)\\
 &\qquad=-\frac{\eta}{12}
 \left(x_1x_2x_3+
 \frac{\hbar\kappa x_3}{(z_1-z_2)^2}+
 \frac{\hbar\kappa x_2}{(z_1-z_3)^2}+
 \frac{\hbar\kappa x_1}{(z_2-z_3)^2}\right).
 \end{split}
\end{equation*}
The output has topological degree one.  The inputs have total degree two,
so $m_3$ has degree $-1$ as required.  The first term survives at
$\hbar=0$; the other terms retain the quantum central contraction.
Applying $\mathcal N$ identifies the parenthesis with
$J(z_1)J(z_2)J(z_3)$ in the chosen radial expansion.
\end{example}

The arity-three identity on $(v,v,w)$ is an exact cancellation:
\[
 \underbrace{\tfrac12w-\tfrac14w}_{b_2\text{ insertions}}
 +\underbrace{(-\tfrac16w)-\tfrac1{12}w}_{b_1\text{ insertions into }b_3}
 =0.
\]
It supplies the homotopy for the binary associator $\eta/4$.
For $(v,w,v,w)$ the arity-four identity is
\[
 -\tfrac1{12}w+\tfrac1{24}w-\tfrac1{24}w+\tfrac1{12}w=0.
\]
The first three summands insert $b_2$ into $b_3$.  The last inserts
$b_3$ into $b_2$.  The remaining summand is zero.  With field
coefficients, every term multiplies the same fourfold Wick product by
Proposition~\ref{prop:interval-wick-associativity}.  These cancellations
therefore verify the identities on the actual holomorphic expressions.

The quartic operation for this contraction vanishes identically.
The following arity nevertheless requires a new operation.  The next
integrations give
\[
 \begin{split}
 F_3(v,w,w)&=s\left(\frac{t^3}{6}-\frac{t^2}{4}+\frac{t}{12}\right),\\
 F_4(v,w,w,w)&=-s\frac{t^2(1-t)^2}{24},\\
 b_4(v,w,w,w)&=0,\qquad b_5(v,w,w,w,w)=-\frac1{720}w.
 \end{split}
\]
Here $\int_0^1t^2(1-t)^2dt=1/30$.  On $(v,w,w,v,w)$, the
three $b_3$-into-$b_3$ terms sum to
\[
 -\frac1{72}w-\frac1{144}w+\frac1{72}w=-\frac1{144}w.
\]
The two nonzero insertions of $b_1$ into $b_5$ give
$w/180+w/720=w/144$.  Their coefficients follow from the five
possible positions of the function input, computed below.  Their cancellation proves the arity-five
identity on an input where setting $b_5=0$ would fail.

\section{Every arity from one polynomial integral}

Let the Bernoulli numbers be defined by
\[
 \frac{u}{e^u-1}=\sum_{k\ge0}B_k\frac{u^k}{k!},
 \qquad B_1=-\frac12.
\]
For $k\ge1$, put $c_k=(-1)^k B_k/k!$.

\begin{proposition}[Complete interval operations]
\label{prop:interval-bernoulli-operations}
For $k\ge2$ and $0\le j\le k$,
\begin{equation*}
 b_{k+1}(w^{\otimes j},v,w^{\otimes(k-j)})
 =(-1)^j\binom{k}{j}c_k w.
\end{equation*}
All other operations of arity at least three on $v,w$ vanish.
These formulas, the binary table, and strict unitality determine every
operation.  For disjoint field supports each value is multiplied by the
full ordered Wick product of its coefficients.
\end{proposition}
\begin{proof}
Every nonzero higher transfer tree has only one function input.
Indeed, every nonroot internal edge applies $h$, so its nonzero output
is a function vanishing at the endpoints.  A vertex cannot multiply
two such outputs and then survive another $h$.  Thus the tree must
be a comb, with one leaf attached at each successive vertex.  Every
attached leaf before the root must be $dt$, since $h$ kills functions.
If the last leaf is a function, the root product is a function
vanishing at both endpoints, which $p$ kills.  An all-$dt$ tree vanishes
at its first vertex.  This proves the asserted restriction on inputs.

Define $T(q)=h(q\,dt)$ on polynomials.  For the word
$(v,w,\ldots,w)$ there is one comb, and its root coefficient is
\[
 c_k=\int_0^1(-T)^{k-1}t\,dt.
\]
To evaluate all these integrals, set
$G(u,t)=\sum_{a\ge0}u^a(-T)^a t$ and
$M(u)=\int_0^1G(u,t)dt$.  The equation $G=t-uTG$ implies
\[
 \partial_tG=1-uG+uM,\qquad G(u,0)=0,\qquad G(u,1)=1.
\]
Solving this formal differential equation gives
\[
 G(u,t)=\frac{1-e^{-ut}}{1-e^{-u}},\qquad
 M(u)=\frac1{1-e^{-u}}-\frac1u
 =\sum_{k\ge1}\frac{(-1)^kB_k}{k!}u^{k-1}.
\]
The apparent poles at $u=0$ cancel, so these are formal power series.
Finally, placing $j$ copies of $w$ before $v$ gives
$\binom{k}{j}$ combs.  Each left attachment contributes a minus sign
from $D_2(s\,dt,sf)=-s(fdt)$.  Right attachments contribute no sign.
Their common integral is $c_k$, proving the formula.
\end{proof}

The scalar formulas follow from the interval integrals, while the
field coefficients enter through the proved Wick comparison.
At $\kappa=0$ or $\hbar=0$, the contraction terms disappear while the
topological higher products remain.  Thus this higher arity is not a
measurement of quantum loop order.

\section{The beta--gamma fields and the curve operation}

For the bosonic beta--gamma system use the base $\C\llbracket\hbar\rrbracket$
and replace the Heisenberg symbols by
$x_{i,r}^{\beta},x_{i,r}^{\gamma}$.  Its elementary contractions are
\[
 C_{ij}^{\beta\gamma}=\frac{\hbar}{z_i-z_j},\qquad
 C_{ij}^{\gamma\beta}=-\frac{\hbar}{z_i-z_j},\qquad
 C_{ij}^{\beta\beta}=C_{ij}^{\gamma\gamma}=0.
\]
Jet contractions are their corresponding holomorphic derivatives.
These coefficients are symmetric under simultaneous exchange of the
point and field labels, as bosonic locality requires.

For completeness, take modes
$\beta(z)=\sum_m\beta_mz^{-m-1}$ and
$\gamma(z)=\sum_n\gamma_nz^{-n}$ with
$[\beta_m,\gamma_n]=\hbar\delta_{m+n,0}$.
The vacuum is killed by $\beta_m$ for $m\ge0$ and by $\gamma_n$
for $n\ge1$.  Explicitly, the vacuum module is
$\C[q_m\ (m\ge1),p_n\ (n\ge0)]\llbracket\hbar\rrbracket$, with
\[
 \beta_{-m}=q_m\ (m\ge1),\qquad
 \gamma_{-n}=p_n\ (n\ge0),\qquad
 \beta_m=\hbar\partial_{p_m}\ (m\ge0),\qquad
 \gamma_n=-\hbar\partial_{q_n}\ (n\ge1).
\]
Summing the two geometric series gives the displayed contractions.
The proof of Proposition~\ref{prop:interval-wick-associativity}
therefore applies without change to the enlarged symbol set.

In particular, writing $\beta_i,\gamma_i$ for its undifferentiated
symbols, the ternary operation is
\[
 m_3(e_1\beta_1,\eta\gamma_2,\eta\beta_3)
 =-\frac{\eta}{12}\left(
 \beta_1\gamma_2\beta_3+
 \frac{\hbar\beta_3}{z_1-z_2}
 -\frac{\hbar\beta_1}{z_2-z_3}\right).
\]
The same scalar higher operations therefore act on both a double-pole
current and a simple-pole first-order system, with their different
holomorphic contractions intact.

There is a further operation when insertion points collide.  In the
Heisenberg vertex algebra, the regular coefficient is
$a_{(-1)}b=\operatorname{Res}_{z=0}z^{-1}Y(a,z)b\,dz$.
This residue extracts a Laurent coefficient and has contour
normalization $1/(2\pi i)$.  For three currents, direct normal ordering
and Taylor expansion give
\[
 ((J_{(-1)}J)_{(-1)}J)-J_{(-1)}(J_{(-1)}J)
 =\hbar\kappa\,\partial^2J.
\]
Indeed, $:JJ:(z)J(w)$ contains
$2\hbar\kappa J(z)/(z-w)^2$, whose regular coefficient contributes
$\hbar\kappa\partial^2J$.  The reversed iterated regular product
has no such term.  Thus this regular product cannot be the binary
multiplication of a degree-zero ordinary minimal $A_\infty$ algebra.
The nonzero result also exhibits the precise information lost by
replacing the separated Wick product with coincident commutative symbols.

\clearpage
\section{Topological resolution and boundary observables}

The interval construction has a strict comparison to the original
multilocal field system.  The map
\[
 W_S\longrightarrow C(S),\qquad f\longmapsto(e_0+e_1)f
\]
is a strict $A_\infty$ quasi-isomorphism.  Its image is the constant
cochain, the differential kills it, and every higher operation on it
vanishes by strict unitality in the interval factor.  For $a=0,1$, evaluation $\epsilon_a:C(S)\to W_S$ is also a strict
quasi-isomorphism: higher outputs have edge degree, and evaluation kills them.
Writing $\iota(f)=(e_0+e_1)f$, the composite $\epsilon_a\iota$ is the
identity.  The other composite has the explicit chain homotopy
\[
 dk_a+k_ad=1-\iota\epsilon_a,\qquad
 k_0(\eta)=e_1,\qquad k_1(\eta)=-e_0,\qquad k_a(e_b)=0.
\]
For example, on $e_0,e_1,\eta$, the first identity with $a=0$
reads $-e_1,e_1,\eta$ on both sides.  This explains
how the same free fields admit a strict model and the explicit nonzero
higher operations calculated above.

For a physical realization of the algebraic specialization
$\kappa=k\in\C$, use the split holomorphic BF family of
Chapter~\ref{chap:zero-level-bf} with bulk parameter $\ell=-k$.
Its fields on $\C\times\R_{\ge0}$ have separate partial-form and
holomorphic-one-form components $(\alpha,\beta)$; write
$d'=\bar\partial+d_t$ for the partial de Rham differential.
They carry a fixed
nondegenerate cross pairing and differential
\[
 Q_\ell(\alpha,\beta)=(d'\alpha,d'\beta+\ell\partial\alpha).
\]
The chiral boundary condition is $\iota^*\alpha=0$.  The integrated
pairing requires at least one compactly supported argument, and the
integrated action is defined on compact fields; local densities are
defined on unrestricted smooth fields.
For $k\ne0$, the field change
$\alpha+\ell^{-1}\beta$ identifies this theory with abelian
Chern--Simons theory whose symmetric coefficient pairing is
$\langle e,e\rangle=\ell=-k$.

The boundary quantum observable complex has the positive unary
operator and contraction
\[
 D=\bar\partial+\hbar\Delta_{\mu_\ell},\qquad
 \mu_\ell(u,v)=-\ell\int_\C^{\mathrm{res}}u\wedge\partial v,
 \qquad \int^{\mathrm{res}}=\frac1{2\pi i}\int.
\]
Choose a smooth radial compact cutoff $\chi$ equal to one near the
origin.  Its current is fixed by
$q_{n+1}=[-\bar\partial\chi/z^{n+1}]$, $n\ge0$.
For separated evaluation forms $u_a,u_b$ and their representatives
$e_a,e_b$ on a common outer radial annulus, the compact chain
calculation gives
\[
 [u_au_b]=[e_ae_b]-\frac{\hbar\ell}{(a-b)^2}\mathbf1
                  =[e_ae_b]+\frac{\hbar k}{(a-b)^2}\mathbf1.
\]
Thus it is the extracted level $-\ell=k$ that matches the algebraic
parameter of this chapter.  The Cauchy conjugation and moment
calculation in Chapter~\ref{chap:zero-level-bf} transport every
separated product to its Wick formula and recover the complete mode action
\[
 J_{-m}=q_m,\qquad J_m=\hbar k m\partial_{q_m}\quad(m>0),\qquad J_0=0,
 \qquad J(z)J(w)\sim\frac{\hbar k}{(z-w)^2}.
\]
The compact collar map intertwines the actual quantum differentials
and all separated products.  Consequently the represented fields
have precisely the normalization used by the interval transfer.
The extraction first uses polynomial $\hbar$ and the finite-weight
vacuum space.  Formal completion is applied afterwards to the
extracted modes and Wick products, coefficient by coefficient; on
arbitrary completed states the field target is the inverse limit
of the Laurent-series spaces modulo $\hbar^M$.  This does not identify
full disk cohomology with its finite-weight vacuum or interchange
cohomology and completion of the full observable complex.

At $k=0$, the split field pairing remains nondegenerate and the
boundary current becomes commutative, with all polynomial jet fields
retained.  The field change to the nonzero de Rham presentation
contains $\ell^{-1}$ and has no zero specialization.  In particular,
the zero symmetric pairing on a single Chern--Simons coefficient
space supplies no inverse BV pairing or propagator.  The zero-level
physical realization uses the two separate BF field components.

The compact quantum contraction of
Theorem~\ref{thm:bf-compact-quantum-contraction} supplies a chain map
from the finite current part of this interval model to the full
boundary observable complex.  Fix its finite list of points and divided
jets, and set $\kappa=-k$ in that theorem.  Its normal polynomial algebra
$A=\C[y_1,\ldots,y_N][\hbar]$ represents the common-annulus currents.
An undivided jet of order $n_i$ is represented by $n_i!y_i$.
Insertion groups remain at distinct holomorphic points.

\begin{proposition}[The supported interval chain map]
\label{prop:interval-supported-current-chain-map}
Let $\iota_U:A\to\mathcal B_{-k}(U)$ and
$L_U:\mathcal B_{-k}(U)\to A$ be the compact inclusion and moment
projection of Proposition~\ref{prop:bf-compact-full-boundary-split}.
For $a=0,1$, define
\[
 \Phi_a:C_I\otimes A\longrightarrow\mathcal B_{-k}(U),
 \qquad \Phi_a(\alpha\otimes F)=\epsilon_a(\alpha)\iota_U(F),
 \qquad
 \Psi(b)=(e_0+e_1)\otimes L_U(b).
\]
These are chain maps for the interval differential and the actual
quantum differential, respectively.  They satisfy
\[
 \Phi_a\Psi=\iota_UL_U,\qquad
 1-\Psi\Phi_a=d(k_a\otimes1)+(k_a\otimes1)d
               \quad\hbox{on }C_I\otimes A.
\]
Their induced map identifies $H^\bullet(C_I\otimes A)=A$ with
the annular polynomial subspace of $H^0(\mathcal B_{-k}(U))$.
For every tuple of polynomial currents at distinct points, the
comparison of its supported product with the annular Wick expression
has the compact primitive $\rho H F$ in
\eqref{eq:bf-compact-wick-chain}.
\end{proposition}
\begin{proof}
Endpoint evaluation kills $dC_I$ and $\iota_U(F)$ is a quantum cycle.
The map $L_U$ intertwines the quantum differential with zero.
These statements prove the chain-map equations.
The equality $L_U\iota_U=1$ and the displayed interval contraction
prove both composite identities.  The split inclusion proves the
cohomology assertion.  The product statement is
\eqref{eq:bf-compact-wick-chain}, with $-\hbar\kappa=\hbar k$.
\end{proof}

One can retain the interval factor during the compact contraction.
On $C_I\otimes Q$, use the total differential
$d(\alpha\otimes q)=d_I\alpha\otimes q+
(-1)^{|\alpha|}\alpha\otimes D_Qq$.  Then
\[
 \mathscr H(\alpha\otimes q)=(-1)^{|\alpha|}\alpha\otimes Hq,
 \qquad
 d\mathscr H+\mathscr Hd=1-(1\otimes\iota\Pi).
\]
The two terms containing $d_I\alpha$ cancel by their opposite signs.
The remaining terms are exactly the compact quantum contraction.
Thus this extension keeps the interval differential distinct from
the quantum differential while intertwining both.

The map $\Phi_a$ is a map of complexes.  The product comparison uses
compact homotopies because common-annulus representatives overlap.
In particular, the second-order quantum differential does not define
an ordinary differential graded multiplication on that full target.
An ordered higher morphism requires compatible Taylor maps on these
same supported complexes.  A collision comparison also requires
operations on collision chains and their boundary strata.
Neither construction follows from the linear chain map alone.

Ordered defect operations and symmetric perturbative brackets likewise
require distinct insertion rules: the former retain the order of
inputs along the defect, while the latter sum over permutations with
their graded signs.  The Wick field comparison
preserves ordered products at distinct points.  Its supported realization
has the compact chain homotopies above.  A comparison with renormalized
configuration integrals must construct compatible Taylor maps on the
same complex and verify their equations at every arity.  Neither the Gaussian contractions nor the interval identities
alone supply the quantum master equation for an interacting theory.

\input{chapters/current-taylor-maps.tex}
